% paper/sections/10_1b_dccd_dissipation.tex
%
% §10.1.5  DCCD フィルタ：分散・散逸特性と選択的高周波減衰の評価
%          実験スクリプト: experiments/dccd_comparison.py
%          結果データ:    results/dccd_comparison/

\noindent\textbf{目的：}
DCCD の 10 次選択フィルタ（§~\ref{sec:advection_dccd_design}）が，
不連続プロファイルに対して CCD の過大な全変動（TV）を大幅に低減しつつ，
滑らかなプロファイルでは CCD と同等の高精度を維持することを検証する．

\noindent\textbf{評価対象：}
DCCD の分散・散逸特性．CLS 法の界面 $\psi$（準不連続プロファイル）に対する
DCCD の選択的高波数減衰の有効性を間接的に評価する．

\noindent\textbf{手順とパラメタ：}
1次元線形移流方程式
\begin{equation}
  \frac{\partial u}{\partial t} + c\,\frac{\partial u}{\partial x} = 0,
  \qquad x \in [0,1),\quad c = 1,\quad T = 1
  \label{eq:linear_advection_dccd}
\end{equation}
を以下の条件で解く：
\begin{itemize}[noitemsep]
  \item 格子 $N=256$，周期境界，CFL $=0.4$，時間積分 RK4，1 周期（$T=1$）
  \item 初期条件 3 種：Square（不連続矩形パルス），Triangle（$C^0$ 三角波），Smooth tanh
  \item 比較スキーム：O2（2次中心差分），O4（4次），CCD（6次），
        DCCD（CCD $+$ 10次選択フィルタ，$\alpha_f=0.4$），WENO5
  \item 評価指標：$L_2 = \sqrt{\mathrm{mean}(u-u_\mathrm{exact})^2}$，
        $\mathrm{TV} = \sum_i |u_{i+1}-u_i|$
\end{itemize}

\noindent\textbf{結果：}

\begin{table}[htbp]
\centering
\caption{1周期後の $L_2$ 誤差（$N=256$，CFL$=0.4$，太字：各初期条件での最良値）}
\label{tab:dccd_l2_bench}
\begin{tabular}{lrrrrr}
\toprule
初期条件 & O2 & O4 & CCD & DCCD & WENO5 \\
\midrule
Square（不連続） & $1.41\times10^{-1}$ & $9.08\times10^{-2}$ & $4.80\times10^{-2}$ & $\mathbf{4.23\times10^{-2}}$ & $6.35\times10^{-2}$ \\
Triangle（$C^0$） & $2.41\times10^{-2}$ & $5.84\times10^{-3}$ & $\mathbf{1.37\times10^{-3}}$ & $1.41\times10^{-3}$ & $3.90\times10^{-3}$ \\
Smooth tanh & $3.75\times10^{-2}$ & $2.46\times10^{-3}$ & $\mathbf{2.57\times10^{-5}}$ & $3.75\times10^{-5}$ & $9.66\times10^{-4}$ \\
\bottomrule
\end{tabular}
\end{table}

\begin{table}[htbp]
\centering
\caption{1周期後の全変動（TV）．
  Exact 値は初期条件の TV（Square: 2.000, Triangle: 1.969, tanh: 2.000）．
  太字：各初期条件での最良値（TV $\to$ exact に最近）．}
\label{tab:dccd_tv_bench}
\begin{tabular}{lrrrrr}
\toprule
初期条件 & O2 & O4 & CCD & DCCD & WENO5 \\
\midrule
Square（不連続） & $18.40$ & $17.17$ & $10.83$ & $3.15$ & $\mathbf{2.00}$ \\
Triangle（$C^0$） & $2.60$ & $2.37$ & $2.09$ & $1.97$ & $\mathbf{1.91}$ \\
Smooth tanh & $2.82$ & $2.02$ & $2.00$ & $2.00$ & $\mathbf{2.00}$ \\
\bottomrule
\end{tabular}
\end{table}

\noindent\textbf{結論：}
\begin{itemize}[noitemsep]
  \item \textbf{不連続（Square）：} DCCD は CCD に比べ TV を $10.83 \to 3.15$ に低減（約 $1/3$）し，
        $L_2$ 誤差も最良（$4.23\times10^{-2}$）であった．
  \item \textbf{平滑（Smooth tanh）：} DCCD と CCD の $L_2$ 誤差は同オーダー
        （$3.75\times10^{-5}$ vs $2.57\times10^{-5}$，1.5 倍以内），TV も同値（$2.000$）．
        フィルタが高波数成分のみを対象とし，滑らかな成分を保護することを確認した．
  \item \textbf{CLS 法への含意：} CLS の界面 $\psi$ は Smooth tanh に相当するため，
        DCCD は CCD と同等の精度で界面を移流しつつ振動を抑制できる．
        §~\ref{sec:interface_advection_bench}（Zalesak・単一渦）の低体積誤差の理論的根拠となる．
\end{itemize}
